\chapter{控制编程与递归自主性的边界}
\begin{chapteroverviewbox}
这一章，我们围绕“控制编程与递归自主性的边界”做一件克制的事：把直觉压缩为状态、关系、时间尺度和可检验后果。它在全书中的任务，是说明多回路稳定性和有界自修改。
\end{chapteroverviewbox}
\ChapterArt{\number\value{chapter}}

\section*{理论来路与依赖接口}
本章的直接思想来源是原文关于控制工程、培育工程和心理工程的生命周期衍生讨论。我们采用原文较晚形成的定义来整理较早讨论，不把对话中的试探性比喻与最终模型并列成互相竞争的“定律”。已有理论锚点主要是鲁棒控制、时延系统、稳定裕度与权限治理（Kalman、Ashby 与混杂系统理论）。

本章使用上一章“多回路稳定性与相变控制”已经得到的对象与边界；本章产出将供下一章“为什么 Agent 需要“培育”而不仅是训练”使用。若后文术语偶尔出现，只作为路线提示，不参与本章推导。

\section*{从哪里出发}
我们已经拥有上一章给出的概念，现在只向前走一步。我会先把对象放进闭环，再讨论它的数学形式，最后检查工程边界。读者不必预先相信本章术语；只需暂时接受一个方法论约定：智能结构的意义来自它在现实反馈中的作用，而不是来自名称本身。

令系统在观察窗口 $[t,t+T]$ 内的状态轨迹为 $x_{[t,t+T]}$，环境扰动为 $d$，行动为 $u$。本章讨论的对象若确实有效，就应改变条件分布 $p(x_{t+T}\mid x_t,u,d)$，或改变达到同一结果所需的代价。这个起点把讨论锚定在鲁棒控制、时延系统、稳定裕度与权限治理上。
\section{哪些变量可修改}
\begin{narration}
现在把镜头移到“哪些变量可修改”。我们只沿本章已经取得的概念向前一步：先确定它在闭环里的角色，再检查它的尺度与失败边界。
\end{narration}

本节把“哪些变量可修改”落实为由治理层明确允许主体调整、明确禁止调整和需要外部授权调整的变量集合。这个定义不是说“任何相似现象都算”，而是要求我们同时给出对象域、允许的干预以及观察窗口。对于主题“控制编程与递归自主性的边界”而言，它承担的是多回路稳定性和有界自修改中的一个可辨识环节。

把这一关系写成最小数学骨架，可得

\begin{equation}\mathcal X=\mathcal X_{free}\cup\mathcal X_{guarded}\cup\mathcal X_{forbidden}.\end{equation}

式中的符号只承担结构角色，具体参数仍需由任务和身体数据辨识。我们首先检查可观测量，再检查控制量，最后检查比它更慢的约束。这样的顺序来自鲁棒控制、时延系统、稳定裕度与权限治理，但本书对哪些变量可修改的命名和组合仍是需要实验验证的建模提案。

这一小节的反例同样重要：没有禁区、预算与回滚点的递归修改会破坏身份和安全不变量。因此评价它不能只看一次任务结果，而要比较干预前后的轨迹、恢复时间、维护成本和风险。如果改变该机制而所有允许观察下的分布都不变，那么它至少在当前实验族中是冗余描述。

\begin{thought}
对“哪些变量可修改”做一次消融：冻结它、删除它，或只允许它以相邻时间尺度交换残差。哪些现象立即消失，哪些只在长期出现？若答案无法区分三种干预，我们还没有找到正确的状态变量。
\end{thought}

最后，把结论交还现实：本节只建立功能层关系，不由此推出特定脑区实现、主观意识或宇宙目的。可测状态、干预后果和明确失败条件，才是从原文直觉走向可检验理论的桥。

\section{哪些 Goal 可以自主产生}
\begin{narration}
现在把镜头移到“哪些 Goal 可以自主产生”。我们只沿本章已经取得的概念向前一步：先确定它在闭环里的角色，再检查它的尺度与失败边界。
\end{narration}

本节把“哪些 Goal 可以自主产生”落实为由治理层明确允许主体调整、明确禁止调整和需要外部授权调整的变量集合。这个定义不是说“任何相似现象都算”，而是要求我们同时给出对象域、允许的干预以及观察窗口。对于主题“控制编程与递归自主性的边界”而言，它承担的是多回路稳定性和有界自修改中的一个可辨识环节。

把这一关系写成最小数学骨架，可得

\begin{equation}\mathcal X=\mathcal X_{free}\cup\mathcal X_{guarded}\cup\mathcal X_{forbidden}.\end{equation}

式中的符号只承担结构角色，具体参数仍需由任务和身体数据辨识。我们首先检查可观测量，再检查控制量，最后检查比它更慢的约束。这样的顺序来自鲁棒控制、时延系统、稳定裕度与权限治理，但本书对哪些 Goal 可以自主产生的命名和组合仍是需要实验验证的建模提案。

这一小节的反例同样重要：没有禁区、预算与回滚点的递归修改会破坏身份和安全不变量。因此评价它不能只看一次任务结果，而要比较干预前后的轨迹、恢复时间、维护成本和风险。如果改变该机制而所有允许观察下的分布都不变，那么它至少在当前实验族中是冗余描述。

\begin{thought}
对“哪些 Goal 可以自主产生”做一次消融：冻结它、删除它，或只允许它以相邻时间尺度交换残差。哪些现象立即消失，哪些只在长期出现？若答案无法区分三种干预，我们还没有找到正确的状态变量。
\end{thought}

最后，把结论交还现实：本节只建立功能层关系，不由此推出特定脑区实现、主观意识或宇宙目的。可测状态、干预后果和明确失败条件，才是从原文直觉走向可检验理论的桥。

\section{哪些权限可以成长}
\begin{narration}
现在把镜头移到“哪些权限可以成长”。我们只沿本章已经取得的概念向前一步：先确定它在闭环里的角色，再检查它的尺度与失败边界。
\end{narration}

本节把“哪些权限可以成长”落实为由治理层明确允许主体调整、明确禁止调整和需要外部授权调整的变量集合。这个定义不是说“任何相似现象都算”，而是要求我们同时给出对象域、允许的干预以及观察窗口。对于主题“控制编程与递归自主性的边界”而言，它承担的是多回路稳定性和有界自修改中的一个可辨识环节。

把这一关系写成最小数学骨架，可得

\begin{equation}\mathcal X=\mathcal X_{free}\cup\mathcal X_{guarded}\cup\mathcal X_{forbidden}.\end{equation}

式中的符号只承担结构角色，具体参数仍需由任务和身体数据辨识。我们首先检查可观测量，再检查控制量，最后检查比它更慢的约束。这样的顺序来自鲁棒控制、时延系统、稳定裕度与权限治理，但本书对哪些权限可以成长的命名和组合仍是需要实验验证的建模提案。

这一小节的反例同样重要：没有禁区、预算与回滚点的递归修改会破坏身份和安全不变量。因此评价它不能只看一次任务结果，而要比较干预前后的轨迹、恢复时间、维护成本和风险。如果改变该机制而所有允许观察下的分布都不变，那么它至少在当前实验族中是冗余描述。

\begin{thought}
对“哪些权限可以成长”做一次消融：冻结它、删除它，或只允许它以相邻时间尺度交换残差。哪些现象立即消失，哪些只在长期出现？若答案无法区分三种干预，我们还没有找到正确的状态变量。
\end{thought}

最后，把结论交还现实：本节只建立功能层关系，不由此推出特定脑区实现、主观意识或宇宙目的。可测状态、干预后果和明确失败条件，才是从原文直觉走向可检验理论的桥。

\section{Ψ 的可塑区}
\begin{narration}
现在把镜头移到“Ψ 的可塑区”。我们只沿本章已经取得的概念向前一步：先确定它在闭环里的角色，再检查它的尺度与失败边界。
\end{narration}

本节把“Ψ 的可塑区”落实为由治理层明确允许主体调整、明确禁止调整和需要外部授权调整的变量集合。这个定义不是说“任何相似现象都算”，而是要求我们同时给出对象域、允许的干预以及观察窗口。对于主题“控制编程与递归自主性的边界”而言，它承担的是多回路稳定性和有界自修改中的一个可辨识环节。

把这一关系写成最小数学骨架，可得

\begin{equation}\mathcal X=\mathcal X_{free}\cup\mathcal X_{guarded}\cup\mathcal X_{forbidden}.\end{equation}

式中的符号只承担结构角色，具体参数仍需由任务和身体数据辨识。我们首先检查可观测量，再检查控制量，最后检查比它更慢的约束。这样的顺序来自鲁棒控制、时延系统、稳定裕度与权限治理，但本书对Ψ 的可塑区的命名和组合仍是需要实验验证的建模提案。

这一小节的反例同样重要：没有禁区、预算与回滚点的递归修改会破坏身份和安全不变量。因此评价它不能只看一次任务结果，而要比较干预前后的轨迹、恢复时间、维护成本和风险。如果改变该机制而所有允许观察下的分布都不变，那么它至少在当前实验族中是冗余描述。

\begin{thought}
对“Ψ 的可塑区”做一次消融：冻结它、删除它，或只允许它以相邻时间尺度交换残差。哪些现象立即消失，哪些只在长期出现？若答案无法区分三种干预，我们还没有找到正确的状态变量。
\end{thought}

最后，把结论交还现实：本节只建立功能层关系，不由此推出特定脑区实现、主观意识或宇宙目的。可测状态、干预后果和明确失败条件，才是从原文直觉走向可检验理论的桥。

\section{Ω/DGA 的可塑区}
\begin{narration}
现在把镜头移到“Ω/DGA 的可塑区”。我们只沿本章已经取得的概念向前一步：先确定它在闭环里的角色，再检查它的尺度与失败边界。
\end{narration}

本节把“Ω/DGA 的可塑区”落实为由治理层明确允许主体调整、明确禁止调整和需要外部授权调整的变量集合。这个定义不是说“任何相似现象都算”，而是要求我们同时给出对象域、允许的干预以及观察窗口。对于主题“控制编程与递归自主性的边界”而言，它承担的是多回路稳定性和有界自修改中的一个可辨识环节。

把这一关系写成最小数学骨架，可得

\begin{equation}\mathcal X=\mathcal X_{free}\cup\mathcal X_{guarded}\cup\mathcal X_{forbidden}.\end{equation}

式中的符号只承担结构角色，具体参数仍需由任务和身体数据辨识。我们首先检查可观测量，再检查控制量，最后检查比它更慢的约束。这样的顺序来自鲁棒控制、时延系统、稳定裕度与权限治理，但本书对Ω/DGA 的可塑区的命名和组合仍是需要实验验证的建模提案。

这一小节的反例同样重要：没有禁区、预算与回滚点的递归修改会破坏身份和安全不变量。因此评价它不能只看一次任务结果，而要比较干预前后的轨迹、恢复时间、维护成本和风险。如果改变该机制而所有允许观察下的分布都不变，那么它至少在当前实验族中是冗余描述。

\begin{thought}
对“Ω/DGA 的可塑区”做一次消融：冻结它、删除它，或只允许它以相邻时间尺度交换残差。哪些现象立即消失，哪些只在长期出现？若答案无法区分三种干预，我们还没有找到正确的状态变量。
\end{thought}

最后，把结论交还现实：本节只建立功能层关系，不由此推出特定脑区实现、主观意识或宇宙目的。可测状态、干预后果和明确失败条件，才是从原文直觉走向可检验理论的桥。

\section{Checkpoint Governance}
\begin{narration}
现在把镜头移到“Checkpoint Governance”。我们只沿本章已经取得的概念向前一步：先确定它在闭环里的角色，再检查它的尺度与失败边界。
\end{narration}

本节把“Checkpoint Governance”落实为由治理层明确允许主体调整、明确禁止调整和需要外部授权调整的变量集合。这个定义不是说“任何相似现象都算”，而是要求我们同时给出对象域、允许的干预以及观察窗口。对于主题“控制编程与递归自主性的边界”而言，它承担的是多回路稳定性和有界自修改中的一个可辨识环节。

把这一关系写成最小数学骨架，可得

\begin{equation}\mathcal X=\mathcal X_{free}\cup\mathcal X_{guarded}\cup\mathcal X_{forbidden}.\end{equation}

式中的符号只承担结构角色，具体参数仍需由任务和身体数据辨识。我们首先检查可观测量，再检查控制量，最后检查比它更慢的约束。这样的顺序来自鲁棒控制、时延系统、稳定裕度与权限治理，但本书对Checkpoint Governance的命名和组合仍是需要实验验证的建模提案。

这一小节的反例同样重要：没有禁区、预算与回滚点的递归修改会破坏身份和安全不变量。因此评价它不能只看一次任务结果，而要比较干预前后的轨迹、恢复时间、维护成本和风险。如果改变该机制而所有允许观察下的分布都不变，那么它至少在当前实验族中是冗余描述。

\begin{thought}
对“Checkpoint Governance”做一次消融：冻结它、删除它，或只允许它以相邻时间尺度交换残差。哪些现象立即消失，哪些只在长期出现？若答案无法区分三种干预，我们还没有找到正确的状态变量。
\end{thought}

最后，把结论交还现实：本节只建立功能层关系，不由此推出特定脑区实现、主观意识或宇宙目的。可测状态、干预后果和明确失败条件，才是从原文直觉走向可检验理论的桥。

\section{Self-Modifiable Set}
\begin{narration}
现在把镜头移到“Self-Modifiable Set”。我们只沿本章已经取得的概念向前一步：先确定它在闭环里的角色，再检查它的尺度与失败边界。
\end{narration}

本节把“Self-Modifiable Set”落实为由治理层明确允许主体调整、明确禁止调整和需要外部授权调整的变量集合。这个定义不是说“任何相似现象都算”，而是要求我们同时给出对象域、允许的干预以及观察窗口。对于主题“控制编程与递归自主性的边界”而言，它承担的是多回路稳定性和有界自修改中的一个可辨识环节。

把这一关系写成最小数学骨架，可得

\begin{equation}\mathcal X=\mathcal X_{free}\cup\mathcal X_{guarded}\cup\mathcal X_{forbidden}.\end{equation}

式中的符号只承担结构角色，具体参数仍需由任务和身体数据辨识。我们首先检查可观测量，再检查控制量，最后检查比它更慢的约束。这样的顺序来自鲁棒控制、时延系统、稳定裕度与权限治理，但本书对Self-Modifiable Set的命名和组合仍是需要实验验证的建模提案。

这一小节的反例同样重要：没有禁区、预算与回滚点的递归修改会破坏身份和安全不变量。因此评价它不能只看一次任务结果，而要比较干预前后的轨迹、恢复时间、维护成本和风险。如果改变该机制而所有允许观察下的分布都不变，那么它至少在当前实验族中是冗余描述。

\begin{thought}
对“Self-Modifiable Set”做一次消融：冻结它、删除它，或只允许它以相邻时间尺度交换残差。哪些现象立即消失，哪些只在长期出现？若答案无法区分三种干预，我们还没有找到正确的状态变量。
\end{thought}

最后，把结论交还现实：本节只建立功能层关系，不由此推出特定脑区实现、主观意识或宇宙目的。可测状态、干预后果和明确失败条件，才是从原文直觉走向可检验理论的桥。

\section{Bounded Self-Modification Envelope}
\begin{narration}
现在把镜头移到“Bounded Self-Modification Envelope”。我们只沿本章已经取得的概念向前一步：先确定它在闭环里的角色，再检查它的尺度与失败边界。
\end{narration}

本节把“Bounded Self-Modification Envelope”落实为由治理层明确允许主体调整、明确禁止调整和需要外部授权调整的变量集合。这个定义不是说“任何相似现象都算”，而是要求我们同时给出对象域、允许的干预以及观察窗口。对于主题“控制编程与递归自主性的边界”而言，它承担的是多回路稳定性和有界自修改中的一个可辨识环节。

把这一关系写成最小数学骨架，可得

\begin{equation}\mathcal X=\mathcal X_{free}\cup\mathcal X_{guarded}\cup\mathcal X_{forbidden}.\end{equation}

式中的符号只承担结构角色，具体参数仍需由任务和身体数据辨识。我们首先检查可观测量，再检查控制量，最后检查比它更慢的约束。这样的顺序来自鲁棒控制、时延系统、稳定裕度与权限治理，但本书对Bounded Self-Modification Envelope的命名和组合仍是需要实验验证的建模提案。

这一小节的反例同样重要：没有禁区、预算与回滚点的递归修改会破坏身份和安全不变量。因此评价它不能只看一次任务结果，而要比较干预前后的轨迹、恢复时间、维护成本和风险。如果改变该机制而所有允许观察下的分布都不变，那么它至少在当前实验族中是冗余描述。

\begin{thought}
对“Bounded Self-Modification Envelope”做一次消融：冻结它、删除它，或只允许它以相邻时间尺度交换残差。哪些现象立即消失，哪些只在长期出现？若答案无法区分三种干预，我们还没有找到正确的状态变量。
\end{thought}

最后，把结论交还现实：本节只建立功能层关系，不由此推出特定脑区实现、主观意识或宇宙目的。可测状态、干预后果和明确失败条件，才是从原文直觉走向可检验理论的桥。

\section*{本章收束：把名词还原为闭环}
现在回头看“控制编程与递归自主性的边界”，最重要的不是记住缩写，而是说清本章对象怎样改变系统的状态、代价或可恢复性。下一章“为什么 Agent 需要“培育”而不仅是训练”只使用这里已经给出的结果，不反过来为本章提供前提。

\begin{bookproposition}
若一个候选机制既不改变可观测轨迹分布，也不改变资源、风险或恢复代价，那么在当前观察尺度上，它不是一个可辨识的功能结构。
\end{bookproposition}
\begin{proofidea}
功能等价由可干预后果定义。若所有允许干预下的输出分布与代价均不变，则该机制相对于当前实验族不可辨识；要么扩大观察尺度，要么删除这一冗余描述。
\end{proofidea}

\begin{readercheck}
请尝试回答：本章对象的状态是什么？它的最快和最慢时间常数是什么？什么观测能证伪它？它失败时，残差应向哪一层升级？
\end{readercheck}
\cleardoublepage
